% Unit 101 terjemahan Bahasa Indonesia dari chapter7.tex baris 2306--2471.
% Sumber: Methods of Algebra, Volume 2, commit
% 9a5803ff77dd3257484cb177f851a73770a59dd3 (CC BY 4.0).
% stable-unit-id: o014.aljabr2.chapter7.galois-descent
% Cakupan: Bagian 7.10 lengkap, ``Penerapan: Penurunan Galois'';
% berhenti tepat sebelum Bagian 7.11 pada baris 2472.

% segment-id: o014.li2.chapter7.galois-descent.h001
\section{Penerapan: Penurunan Galois}\label{sec:descent-Galois}
\hypertarget{o014.aljabr2.chapter7.galois-descent}{ }
% segment-id: o014.li2.chapter7.galois-descent.p001
Dalam pembahasan berikut, $L|K$ adalah perluasan Galois medan. Ini merupakan
kasus khusus yang berguna dari penurunan datar dalam
\S\ref{sec:descent-modules}, dan disebut \emph{penurunan Galois}. Strategi
kita ialah mendeskripsikan komonad yang bersesuaian secara langsung. Alat yang
diperlukan hanya bentuk penurunan datar dalam Teorema~\ref{prop:flat-descent}
beserta persiapannya, tanpa pembahasan lebih mendalam pada paruh kedua
\S\ref{sec:descent-modules}.

% segment-id: o014.li2.chapter7.galois-descent.p002
Mulai sekarang, andaikan $K$ dan $L$ medan serta $L|K$ perluasan Galois.
Modul-modul sebelumnya kini menjadi ruang vektor. Tuliskan
% segment-id: o014.li2.chapter7.galois-descent.q001
\[ \Gamma := \Gal(L|K), \quad \text{dilengkapi topologi Krull}. \]
Latar belakang topologi Krull diberikan dalam
\cite[\S 4.10, \S 9.2]{Li1}; bila $L|K$ hingga, ini hanyalah topologi
diskret. Sesungguhnya, $\Gamma$ adalah grup profinit seperti yang
diperkenalkan dalam \S\ref{sec:profinite-groups}.

% segment-id: o014.li2.chapter7.galois-descent.p003
Suatu pemetaan $f:\Gamma\to X$ disebut \emph{mulus} jika, untuk setiap
$\sigma\in\Gamma$, terdapat subgrup buka $\Gamma_0$ sedemikian sehingga
$\sigma_0\in\Gamma_0\implies f(\sigma\sigma_0)=f(\sigma)$. Tuliskan
% segment-id: o014.li2.chapter7.galois-descent.q002
\[ \mathrm{Maps}(\Gamma, X)^\infty := \left\{ \text{pemetaan mulus } f: \Gamma \to X \right\}. \]

% segment-id: o014.li2.chapter7.galois-descent.le001
\begin{lemma}\label{prop:smooth-map-finite}
	Jika $f\in\mathrm{Maps}(\Gamma,X)^\infty$, terdapat subgrup normal buka
	$\Gamma'$ sedemikian sehingga $f$ terfaktorkan melalui
	$\Gamma/\Gamma'$. Pemetaan $f$ hanya mengambil sejumlah nilai berhingga.
\end{lemma}
% segment-id: o014.li2.chapter7.galois-descent.pf001
\begin{proof}
	Untuk setiap $\sigma\in\Gamma$, terdapat subgrup buka $\Gamma_0$ yang
	bergantung pada $\sigma$ sedemikian sehingga $f$ konstan pada
	$\sigma\Gamma_0$. Koset-koset $\sigma\Gamma_0$ ini membentuk peliput buka
	bagi $\Gamma$. Karena $\Gamma$ kompak, terdapat subpeliput hingga
	$\sigma_1\Gamma_{0,1},\ldots,\sigma_r\Gamma_{0,r}$. Sekarang ambil subgrup
	normal buka $\Gamma'$ yang terkandung dalam
	$\bigcap_{i=1}^r\Gamma_{0,i}$.
\end{proof}

% segment-id: o014.li2.chapter7.galois-descent.p004
Kita gunakan hasil ini untuk mendeskripsikan secara konkret komonad
$\mathbf{L}:M\mapsto L\dotimes{K}M$ yang ditentukan $L|K$ pada
$\cate{Vect}(L)$.

% segment-id: o014.li2.chapter7.galois-descent.le002
\begin{lemma}\label{prop:Galois-descent-comonad-0}
	Untuk setiap ruang vektor-$L$ $M$, berikan struktur ruang vektor-$L$ pada
	$\mathrm{Maps}(\Gamma,M)^\infty$ sebagai berikut: bila $\ell\in L$, maka
	$(\ell f)(\sigma)=\sigma(\ell)f(\sigma)$ untuk setiap $\sigma\in\Gamma$.
	Terdapat isomorfisme ruang vektor-$L$
	% segment-id: o014.li2.chapter7.galois-descent.q003
	\begin{equation}\label{eqn:comonad-smMap}
		\begin{tikzcd}[row sep=tiny]
			\Phi(M): L \dotimes{K} M \arrow[r, "\sim"] & \mathrm{Maps}(\Gamma, M)^\infty \\
			\ell \otimes m \arrow[mapsto, r] & {\left[ f: \sigma \mapsto \sigma(\ell)m \right]}.
		\end{tikzcd}
	\end{equation}
\end{lemma}
% segment-id: o014.li2.chapter7.galois-descent.pf002
\begin{proof}
	Jelas bahwa $\Phi(M)$ terdefinisi dengan baik
	($\Stab_\Gamma(\ell)$ selalu buka, sehingga $f$ mulus) dan, menurut
	definisi, $L$-linear. Pemetaan ini juga kompatibel dengan sembarang jumlah
	langsung:
	$\Phi(\bigoplus_iM_i)=\bigoplus_i\Phi(M_i)$ (setiap $f$ hanya mengambil
	berhingga banyak nilai, sehingga
	$\mathrm{Maps}(\Gamma,\mathord\cdot)^\infty$ mempertahankan sembarang jumlah
	langsung). Dengan memilih basis $M$, \eqref{eqn:comonad-smMap} cukup
	diperiksa untuk kasus khusus $M=L$,
	% segment-id: o014.li2.chapter7.galois-descent.q004
	\[\begin{tikzcd}[row sep=tiny]
		\Phi(L): L \dotimes{K} L \arrow[r, "\sim"] & \mathrm{Maps}(\Gamma, L)^\infty \\
		\ell \otimes \ell' \arrow[mapsto, r] & {[f: \sigma \mapsto \sigma(\ell)\ell' ]}.
	\end{tikzcd}\]

	Untuk sembarang subperluasan Galois hingga $L'|K$ dari $L|K$, tuliskan
	$\Gamma':=\Gal(L|L')$. Kita klaim bahwa pemetaan yang sama seperti di
	atas memberikan
	% segment-id: o014.li2.chapter7.galois-descent.q005
	\begin{equation}\label{eqn:comonad-smMap-aux}
		L' \dotimes{K} L' \rightiso \mathrm{Maps}(\Gamma/\Gamma', L') := \left\{\text{pemetaan } f: \Gamma/\Gamma' \to L' \right\}.
	\end{equation}
	Setelah pernyataan ini diterima, karena
	$\mathrm{Maps}(\Gamma,L)^\infty=
	\bigcup_{L'|K}\mathrm{Maps}(\Gamma/\Gamma',L')$
	(terapkan Lema~\ref{prop:smooth-map-finite}; memperbesar $L'$ sama dengan
	memperkecil $\Gamma'$), mengambil gabungan kedua ruas
	\eqref{eqn:comonad-smMap-aux} atas semua $L'|K$ menunjukkan bahwa
	$\Phi(L)$ adalah isomorfisme.

	Pilih $\alpha$ dengan $L'=K(\alpha)$. Polinomial minimalnya
	$p\in K[X]$ terurai dalam $L'$ sebagai
	$\prod_{i=1}^d(X-\alpha_i)$, dengan $\alpha_1=\alpha$ dan tanpa akar
	berulang. Grup $\Gamma/\Gamma'=\Gal(L'|K)$ berkorespondensi satu-satu
	dengan $\alpha_1,\ldots,\alpha_d$ melalui
	$\sigma\mapsto\sigma(\alpha)$. Tuliskan $\ell\in L'$ sebagai $g(\alpha)$
	dengan $g\in K[X]$. Maka
	% segment-id: o014.li2.chapter7.galois-descent.q006
	\[\begin{tikzcd}[row sep=tiny, column sep=small]
		L' \dotimes{K} L' & \dfrac{K[X]}{(p)} \dotimes{K} L' \arrow[l, "\sim"'] \arrow[r, "\sim"] & \dfrac{L'[X]}{\displaystyle\prod_{i=1}^d (X - \alpha_i)} \arrow[r, "\sim"] & \displaystyle\prod_{i=1}^d \underbracket{\dfrac{L'[X]}{(X - \alpha_i)}}_{\text{diidentifikasi dengan } L'} \arrow[r, "\sim"] & (L')^{\Gamma/\Gamma'} \\
		\ell \otimes \ell' & (g + (p)) \otimes \ell' \arrow[mapsto, l] \arrow[mapsto, r] & g\ell' + (p) \arrow[mapsto, r] & \left( g(\alpha_i) \ell' \right)_{i=1}^d \arrow[mapsto, r] & \left( \sigma_i(\ell) \ell' \right)_{i=1}^d.
	\end{tikzcd}\]
	Ini membuktikan \eqref{eqn:comonad-smMap-aux}.
\end{proof}

% segment-id: o014.li2.chapter7.galois-descent.le003
\begin{lemma}\label{prop:Galois-descent-comonad-1}
	Melalui isomorfisme kanonik dalam
	Lema~\ref{prop:Galois-descent-comonad-0}, komonad
	$(\mathbf{L},\delta,\epsilon)$ dideskripsikan oleh diagram komutatif
	berikut, dengan $M$ sembarang ruang vektor-$L$:
	% segment-id: o014.li2.chapter7.galois-descent.q007
	\[\begin{tikzcd}
		M \arrow[equal, d] & L \dotimes{K} M \arrow[r, "{\delta_M}"] \arrow[d, "\sim"', sloped] \arrow[l, "{\epsilon_M}"'] & L \dotimes{K} L \dotimes{K} M \arrow[d, "\sim", sloped] \\
		M & \mathrm{Maps}(\Gamma, M)^\infty \arrow[r, "{\delta'_M}"] \arrow[l, "{\epsilon'_M}"'] & \mathrm{Maps}\left(\Gamma, (\mathrm{Maps}(\Gamma, M)^\infty )\right)^\infty \\
		f(1) \arrow[phantom, u, "\in" description, sloped] & f \arrow[mapsto, r] \arrow[phantom, u, "\in" description, sloped] \arrow[mapsto, l] & {[\sigma \mapsto [\tau \mapsto f(\tau\sigma) ]]} \arrow[phantom, u, "\in" description, sloped]
	\end{tikzcd}\]
	dengan $1:=1_\Gamma$.
\end{lemma}
% segment-id: o014.li2.chapter7.galois-descent.pf003
\begin{proof}
	Kounit $\epsilon_M:L\dotimes{K}M\to M$ adalah
	$\ell\otimes m\mapsto\ell m$, yang di bawah
	\eqref{eqn:comonad-smMap} jelas bersesuaian dengan $f\mapsto f(1)$.
	Komultiplikasi
	$\delta_M:L\dotimes{K}M\to L\dotimes{K}L\dotimes{K}M$ memerlukan sedikit
	perhitungan. Jika $f\in\mathrm{Maps}(\Gamma,M)^\infty$ bersesuaian dengan
	$\ell\otimes m$, citra yang terakhir di bawah $\delta_M$, yaitu
	$\ell\otimes(1\otimes m)$, bersesuaian dengan pemetaan
	% segment-id: o014.li2.chapter7.galois-descent.q008
	\[ \Gamma \ni \sigma \mapsto \sigma(\ell)  \cdot \underbracket{\left[ \tau \mapsto m  \right]}_{\in \mathrm{Maps}(\Gamma, M)^\infty} = \left[ \tau \mapsto \tau(\sigma(\ell))m \xlongequal{\text{\eqref{eqn:comonad-smMap}}} f(\tau\sigma) \right]. \]
	Selesailah bukti.
\end{proof}

% segment-id: o014.li2.chapter7.galois-descent.d001
\begin{definition}\label{def:semilinear-action}
	\index{aksi semilinear mulus@aksi semilinear mulus (smooth semi-linear action)}
	\index[sym1]{Vect-Gamma-infty@$\cate{Vect}^{\Gamma, \infty}$}
	Misalkan grup $\Gamma=\Gal(L|K)$ beraksi dari kiri pada ruang vektor-$L$
	$M$. Aksi itu disebut \emph{semilinear} bila
	$\sigma(m+m')=\sigma(m)+\sigma(m')$ dan
	$\sigma(tm)=\sigma(t)\sigma(m)$ untuk semua $\sigma\in\Gamma$, $t\in L$,
	dan $m,m'\in M$. Definisikan kategori
	$\cate{Vect}^{\Gamma,\infty}(L)$ sebagai berikut.
	\begin{itemize}
		\item Objek: ruang vektor-$L$ $M$ dengan aksi $\Gamma$ semilinear
		mulus. ``Mulus'' berarti bahwa $\Stab_\Gamma(m)$ adalah subgrup buka
		dari $\Gamma$ untuk setiap $m\in M$.
		\item Morfisme: pemetaan $L$-linear $\varphi$ yang memenuhi
		$\varphi(\sigma m)=\sigma(\varphi(m))$, atau pemetaan linear
		$\Gamma$-ekuivarian.
	\end{itemize}
\end{definition}

% segment-id: o014.li2.chapter7.galois-descent.p005
Kembali ke komonad $\mathbf{L}$ yang ditentukan oleh perluasan Galois $L|K$.
Kategori komodulnya ditulis $\cate{Vect}(L)^{\mathbf{L}}$.

% segment-id: o014.li2.chapter7.galois-descent.le004
\begin{lemma}\label{prop:Galois-descent-aux}
	Dengan notasi di atas, terdapat isomorfisme kategori
	% segment-id: o014.li2.chapter7.galois-descent.q009
	\[ \cate{Vect}(L)^{\mathbf{L}} \simeq \cate{Vect}^{\Gamma, \infty}(L). \]
	Jika $M$ objek $\cate{Vect}^{\Gamma,\infty}(L)$, objek yang bersesuaian
	dalam $\cate{Vect}(L)^{\mathbf{L}}$ adalah $(M,a)$, dengan $a$ dicirikan
	oleh fakta bahwa komposisi
	% segment-id: o014.li2.chapter7.galois-descent.q010
	\[ M \xrightarrow{a} L \dotimes{K} M \xrightarrow[\sim]{\Phi(M)} \mathrm{Maps}(\Gamma, M)^\infty \]
	memetakan $m\in M$ ke pemetaan $[\sigma\mapsto\sigma m]$ dari $\Gamma$ ke
	$M$.
\end{lemma}
% segment-id: o014.li2.chapter7.galois-descent.pf004
\begin{proof}
	Pertama, perhatikan bahwa $\mathrm{Maps}(\Gamma,\mathord\cdot)^\infty$
	adalah funktor: setiap pemetaan $\varphi:M_1\to M_2$ menginduksi
	$\mathrm{Maps}(\Gamma,M_1)^\infty\to
	\mathrm{Maps}(\Gamma,M_2)^\infty$ yang memetakan $f$ ke $\varphi f$.

	Selanjutnya, gunakan Lema~\ref{prop:Galois-descent-comonad-1} untuk
	menulis ulang dan menggabungkan kedua diagram komutatif yang mendefinisikan
	komodul menjadi
	% segment-id: o014.li2.chapter7.galois-descent.q011
	\[\begin{tikzcd}
		\mathrm{Maps}\left( \Gamma, \mathrm{Maps}(\Gamma, M)^\infty \right)^\infty & \mathrm{Maps}(\Gamma, M)^\infty \arrow[l, "{\mathrm{Maps}(\Gamma, a')^\infty }"' inner sep=0.8em] \arrow[r, "{\epsilon'_M}"] & M \\
		\mathrm{Maps}(\Gamma, M)^\infty \arrow[u, "{\delta'_M}"] & M \arrow[l, "{a'}"] \arrow[u, "{a'}"] \arrow[equal, ru] &
	\end{tikzcd}\]
	dengan $a'$ bersesuaian dengan $a$. Memberikan pemetaan $L$-linear
	$a':M\to\mathrm{Maps}(\Gamma,M)^\infty$ ekuivalen dengan memberikan
	pemetaan $\beta:\Gamma\times M\to M$ sedemikian sehingga
	\begin{itemize}
		\item $\beta(\sigma,m+m')=\beta(\sigma,m)+\beta(\sigma,m')$ dan
		$\beta(\sigma,tm)=\sigma(t)\beta(\sigma,m)$, dengan $t\in L$;
		\item untuk setiap $m\in M$ dan $\sigma\in\Gamma$, terdapat subgrup
		buka $\Gamma_0$ sedemikian sehingga
		$\sigma_0\in\Gamma_0\implies
		\beta(\sigma\sigma_0,m)=\beta(\sigma,m)$.
	\end{itemize}
	Korespondensi konkretnya tentu $\beta(\sigma,m)=a'(m)(\sigma)$. Diagram
	komutatif di atas kemudian diterjemahkan menjadi
	% segment-id: o014.li2.chapter7.galois-descent.q012
	\[ \beta(\tau\sigma, m) = \beta(\tau, \beta(\sigma, m)), \quad \beta(1, m) = m. \]

	Pada tingkat objek, syarat-syarat ini persis sama dengan aksi semilinear
	mulus dari $\Gamma$. Pemeriksaan pada tingkat morfisme juga langsung.
\end{proof}

% segment-id: o014.li2.chapter7.galois-descent.le005
\begin{lemma}\label{prop:Galois-descent-action}
	Jika $(M,a)\in\Obj(\cate{Vect}(L)^{\mathbf{L}})$ berasal dari ruang
	vektor-$K$ $N$, objek yang bersesuaian dalam
	$\cate{Vect}^{\Gamma,\infty}(L)$ adalah $L\dotimes{K}N$ dengan aksi
	semilinear mulus
	% segment-id: o014.li2.chapter7.galois-descent.q013
	\[ \sigma (\ell \otimes x) = \sigma(\ell) \otimes x, \quad \ell \in L, \; x \in N, \quad \sigma \in \Gamma. \]
\end{lemma}
% segment-id: o014.li2.chapter7.galois-descent.pf005
\begin{proof}
	Kita mempunyai $M=L\dotimes{K}N$. Menurut
	\eqref{eqn:N-comonad-action}, homomorfisme $a$ memetakan
	$\ell\otimes x$ ke $\ell\otimes1\otimes x$. Jadi, pencirian aksi
	semilinear memberikan
	% segment-id: o014.li2.chapter7.galois-descent.q014
	\[ \sigma (\ell \otimes x) = \Phi(M)(\ell \otimes (1 \otimes x))(\sigma) = \sigma(\ell) (1 \otimes x) = \sigma(\ell) \otimes x, \]
	untuk sembarang $\ell\in L$ dan $x\in N$. Selesailah bukti.
\end{proof}

% segment-id: o014.li2.chapter7.galois-descent.p006
Ambil $(M,a)\in\Obj(\cate{Vect}(L)^{\mathbf{L}})$. Karena
$\Phi(M)(1\otimes m)(\sigma)=\sigma(1)m=m$, deskripsi $a$ dalam
Lema~\ref{prop:Galois-descent-aux} dapat diperluas menjadi diagram komutatif
dalam $\cate{Vect}(K)$,
% segment-id: o014.li2.chapter7.galois-descent.q015
\[\begin{tikzcd}[row sep=large]
	M \arrow[r, "a"] \arrow[rd, "{m \mapsto [\sigma \mapsto \sigma m]}"'] & L \dotimes{K} M \arrow[d, "\Phi(M)" description] & M. \arrow[l, "{1 \otimes m \mapsfrom m}"'] \arrow[ld, "{[m \mapsfrom \sigma] \mapsfrom m}"] \\
	& \mathrm{Maps}(\Gamma, M)^\infty &
\end{tikzcd}\]
Dengan demikian, penyama dalam Teorema~\ref{prop:flat-descent} sama dengan
mengambil invarian-$\Gamma$,
% segment-id: o014.li2.chapter7.galois-descent.q016
\[ M^\Gamma := \left\{ m \in M: \forall \sigma \in \Gamma, \; \sigma m = m \right\}. \]

% segment-id: o014.li2.chapter7.galois-descent.p007
Perhatikan bahwa $M\mapsto M^\Gamma$ memberikan funktor $K$-linear antara
kategori-kategori $K$-linear,
% segment-id: o014.li2.chapter7.galois-descent.q017
\[ (\cdot)^{\Gamma}: \cate{Vect}^{\Gamma, \infty}(L) \to \cate{Vect}(K). \]

% segment-id: o014.li2.chapter7.galois-descent.th001
\begin{theorem}[Penurunan Galois]\label{prop:Galois-descent}
	\index{desen!Galois}
	Misalkan $L|K$ perluasan Galois medan dan
	$\Gamma=\Gal(L|K)$. Maka funktor
	% segment-id: o014.li2.chapter7.galois-descent.q018
	\[\begin{tikzcd}[row sep=tiny]
		\cate{Vect}(K) \arrow[r] & \cate{Vect}^{\Gamma, \infty}(L) \\
		\text{objek } N \arrow[mapsto, r] & L \dotimes{K} N \\
		\text{morfisme } \varphi \arrow[mapsto, r] & \identity \otimes \varphi
	\end{tikzcd}\]
	merupakan ekuivalensi kategori, dengan salah satu funktor kuasi-inversnya
	dapat diambil sebagai
	$(\mathord\cdot)^\Gamma:\cate{Vect}^{\Gamma,\infty}(L)
	\to\cate{Vect}(K)$.
\end{theorem}
% segment-id: o014.li2.chapter7.galois-descent.pf006
\begin{proof}
	Jelas bahwa $L$ adalah aljabar-$K$ yang datar setia.
	Lema~\ref{prop:Galois-descent-aux} dan deskripsi penyama di atas
	menjadikan pernyataan ini kasus khusus Teorema~\ref{prop:flat-descent}
	tentang penurunan datar.
\end{proof}

% segment-id: o014.li2.chapter7.galois-descent.p008
Mari kita tinjau Teorema~\ref{prop:Galois-descent} dari sudut lain. Untuk setiap
$M\in\Obj(\cate{Vect}^{\Gamma,\infty}(L))$, definisikan
% segment-id: o014.li2.chapter7.galois-descent.q019
\begin{equation}\label{eqn:Galois-descent-iotaM}
	\iota_M: L \dotimes{K} (M^{\Gamma}) \to M, \quad
	\ell \otimes x \mapsto \ell x .
\end{equation}
Ini jelas pemetaan linear $\Gamma$-ekuivarian. Terdapat pasangan adjoin
sederhana
% segment-id: o014.li2.chapter7.galois-descent.q020
\begin{equation}\label{eqn:Galois-descent-adjunction}
	\begin{tikzcd}[row sep=tiny]
		L \dotimes{K} (\cdot) : \cate{Vect}(K) \arrow[shift left, r] & \cate{Vect}^{\Gamma, \infty}(L) : (\cdot)^{\Gamma}, \arrow[shift left, l] \\
		\Hom_{\cate{Vect}^{\Gamma, \infty}(L)} (L \dotimes{K} N, M) \arrow[r, "\sim"] & \Hom_{\cate{Vect}(K)}(N, M^{\Gamma}) \\
		\varphi \arrow[mapsto, r] & \varphi|_{1 \otimes N} \\
		\iota_M \circ (\identity_L \otimes \psi) & \psi. \arrow[mapsto, l]
	\end{tikzcd}
\end{equation}
% segment-id: o014.li2.chapter7.galois-descent.l001
\begin{itemize}
	\item Mudah dilihat bahwa unit pasangan adjoin
	$N\to(L\dotimes{K}N)^\Gamma$ memetakan $x$ ke $1\otimes x$. Pemetaan ini
	adalah isomorfisme: ambil basis $B$ dari $N$. Unsur
	$\sum_{b\in B}\ell_b\otimes b$ (jumlah hingga) dari $L\dotimes{K}N$
	bersifat $\Gamma$-invarian jika dan hanya jika setiap $\ell_b$ invarian;
	dengan kata lain,
	$\sum_b\ell_b\otimes b=1\otimes\sum_b\ell_b b\in N$.
	\item Morfisme kounit tepat merupakan $\iota_M$ yang didefinisikan di
	atas.

	Bila $M=L\dotimes{K}N$, butir sebelumnya memberikan
	$L\dotimes{K}(M^\Gamma)=L\dotimes{K}(1\otimes N)\rightiso
	L\dotimes{K}N=M$, sehingga $\iota_M$ adalah isomorfisme dalam kasus ini.
\end{itemize}

% segment-id: o014.li2.chapter7.galois-descent.p009
Setelah Teorema~\ref{prop:Galois-descent} diterima, setiap $M$ berasal dari
suatu $N$, sehingga pembahasan di atas menunjukkan bahwa
$(L\dotimes{K}(\mathord\cdot),(\mathord\cdot)^\Gamma)$ adalah pasangan
ekuivalensi adjoin. Sebaliknya, jika dapat ditunjukkan bahwa $\iota_M$ adalah
isomorfisme untuk setiap $M$, Teorema~\ref{prop:Galois-descent} dapat
dibuktikan kembali. Walaupun bukti melalui penurunan datar mempunyai alur yang
alami, jalannya memang terasa memutar. Latihan akan memberikan metode untuk
membuktikan secara langsung bahwa $\iota_M$ adalah isomorfisme; tekniknya
juga berguna dalam keadaan lain.

% segment-id: o014.li2.chapter7.galois-descent.r001
\begin{remark}[Penurunan Galois bagi aljabar dan modul]\label{rem:Galois-descent-alg}
	Misalkan $N_1,N_2\in\Obj(\cate{Vect}(K))$. Melalui
	\eqref{eqn:change-of-ring-monoidal}, aksi semilinear $\Gamma$ yang alami
	pada $L\dotimes{K}(N_1\dotimes{K}N_2)$ diterjemahkan menjadi
	``aksi diagonal'' pada
	$(L\dotimes{K}N_1)\dotimes{L}(L\dotimes{K}N_2)$, yakni aksi serentak pada
	kedua slot tensor. Memberikan aljabar-$K$ $N$ sama artinya dengan
	memberikan struktur aljabar-$L$ pada $M:=L\dotimes{K}N$ sedemikian sehingga
	operasi aljabarnya kompatibel dengan aksi semilinear $\Gamma$:
	% segment-id: o014.li2.chapter7.galois-descent.q021
	\[ \sigma(1_M) = 1_M , \quad \sigma (xy) = \sigma(x) \sigma(y). \]
	Modul di atas aljabar-$K$ harus dipandang dengan cara yang sama. Gagasan
	ini serupa dengan Catatan~\ref{rem:flat-descent-alg}.
\end{remark}
